% Partie : intuitionistic-logic
% Chapitre : soundness-completeness
% Section : truth-lemma

\documentclass[../../../include/open-logic-section]{subfiles}

\begin{document}

\olfileid[fr]{int}{sc}{tru}

\olsection{Lemme de vérité}

Le lemme suivant relie la vérité d'une !!{formula} au monde~$\sigma$
du modèle canonique à sa !!{derivability} à partir de l'ensemble
premier~$\Delta(\sigma)$.

\begin{lem}\ollabel{lem:truth}
  Si $\Delta$ est premier, alors, pour tout monde~$\sigma$ de
  $\mModel{M(\Delta)}$ et toute !!{formula}~$!A$, on a
  $\mSat{M(\Delta)}{!A}[\sigma]$ si et seulement si
  $\Delta(\sigma) \Proves !A$.
\end{lem}

\begin{proof}
  Raisonnons par induction sur~$!A$, en établissant l'équivalence
  simultanément pour tous les mondes~$\sigma$.
  \begin{enumerate}
    \item \indcase{!A}{\lfalse}{Puisque $\Delta(\sigma)$ est premier,
      il est cohérent, donc $\Delta(\sigma) \Proves/ \indfrm$.
      Par définition, $\mSat/{M(\Delta)}{\indfrm}[\sigma]$.}
    \item \indcase{!A}{p}{Par définition de $\mSat{{}}{}$,
      $\mSat{M(\Delta)}{\indfrm}[\sigma]$ si et seulement si
      $\sigma \in V(p)$, c'est-à-dire si $p\in\Delta(\sigma)$.
      Cela équivaut à $\Delta(\sigma) \Proves \indfrm$, par la
      règle des hypothèses et la clôture de $\Delta(\sigma)$
      par déduction.}
    \item \indcase{!A}{\lnot !B}{Supposons
      $\Delta(\sigma) \Proves \lnot !B$. Pour tout $\sigma'$ tel que
      $R\sigma\sigma'$, on a $\Delta(\sigma)\subseteq\Delta(\sigma')$.
      Si $\Delta(\sigma')$ démontrait aussi $!B$, l'élimination
      de la négation contredirait sa cohérence. Par hypothèse
      d'induction, aucun de ces mondes ne satisfait $!B$, donc
      $\mSat{M(\Delta)}{\indfrm}[\sigma]$.

      Si $\Delta(\sigma) \Proves/ \lnot !B$, alors
      $\Delta(\sigma)\cup\{!B\} \Proves/ \lfalse$, par la règle
      d'introduction de la négation. Soit $n$ l'indice du couple
      $\tuple{!B,\lfalse}$. Par construction,
      $\Delta(\sigma.n)$ contient~$!B$ ; par hypothèse
      d'induction, $\mSat{M(\Delta)}{!B}[\sigma.n]$.
      Comme $R\sigma(\sigma.n)$, la clause de satisfaction de la
      négation donne $\mSat/{M(\Delta)}{\indfrm}[\sigma]$.}
    \item \indcase{!A}{!B \land
      !C}{$\mSat{M(\Delta)}{\indfrm}[\sigma]$ si et seulement si
      $\mSat{M(\Delta)}{!B}[\sigma]$ et
      $\mSat{M(\Delta)}{!C}[\sigma]$. Par hypothèse d'induction,
      $\mSat{M(\Delta)}{!B}[\sigma]$ si et seulement si
      $\Delta(\sigma) \Proves !B$, et de même pour~$!C$.
      Mais $\Delta(\sigma) \Proves !B$ et
      $\Delta(\sigma) \Proves !C$ si et seulement si
      $\Delta(\sigma) \Proves \indfrm$.}
    \item \indcase{!A}{!B \lor
      !C}{$\mSat{M(\Delta)}{\indfrm}[\sigma]$ si et seulement si
      $\mSat{M(\Delta)}{!B}[\sigma]$ ou
      $\mSat{M(\Delta)}{!C}[\sigma]$. Par hypothèse d'induction,
      cela équivaut à $\Delta(\sigma) \Proves !B$ ou
      $\Delta(\sigma) \Proves !C$. Il reste à montrer que cette
      condition équivaut à $\Delta(\sigma) \Proves \indfrm$.
      Le sens direct résulte de l'introduction de la disjonction.
      Réciproquement, si $\Delta(\sigma) \Proves \indfrm$, alors
      $\indfrm \in \Delta(\sigma)$ par clôture par déduction.
      Comme $\Delta(\sigma)$ est premier, il contient $!B$ ou~$!C$.
      Dans l'un et l'autre cas, la formule qu'il contient est
      déductible par la règle des hypothèses.}
    \item \indcase{!A}{!B \lif !C}{D'abord, prouvons la contraposée
      du sens direct. Supposons $\Delta(\sigma) \Proves/ !B
      \lif !C$. Alors $\Delta(\sigma) \cup \{!B\} \Proves/
      !C$. Comme $\tuple{!B, !C}$ est $\tuple{!B_n, !C_n}$ pour
      un certain~$n$, nous avons
      $\Delta(\sigma.n) = (\Delta(\sigma) \cup \{!B\})^*$,
      ainsi que $\Delta(\sigma.n) \Proves !B$ mais
      $\Delta(\sigma.n) \Proves/ !C$. Par hypothèse d'induction,
      $\mSat{M(\Delta)}{!B}[\sigma.n]$ et
      $\mSat/{M(\Delta)}{!C}[\sigma.n]$. Puisque
      $R\sigma(\sigma.n)$, il s'ensuit que
      $\mSat/{M(\Delta)}{\indfrm}[\sigma]$.

      Supposons maintenant $\Delta(\sigma) \Proves !B \lif !C$ et
      soit $\sigma'$ tel que $R\sigma\sigma'$. Puisque
      $\Delta(\sigma) \subseteq \Delta(\sigma')$, si
      $\Delta(\sigma') \Proves !B$, alors
      $\Delta(\sigma') \Proves !C$. Autrement dit, pour tout
      $\sigma'$ tel que $R\sigma\sigma'$, on a soit
      $\Delta(\sigma') \Proves/ !B$, soit
      $\Delta(\sigma') \Proves !C$. Par hypothèse d'induction,
      cela signifie que si $R\sigma\sigma'$, alors
      $\mSat/{M(\Delta)}{!B}[\sigma']$ ou
      $\mSat{M(\Delta)}{!C}[\sigma']$, c'est-à-dire
      $\mSat{M(\Delta)}{\indfrm}[\sigma]$.}
  \end{enumerate}
\end{proof}

\end{document}
